% year: תשפ"ב
% type: מבחן
% semester: א
% moed: ב
% number: 2ב
% points: 9
% categories: func-analysis
% source: exams/מבחנים/תשפב/מועד ב תשפב.pdf

\begin{question}
חשבו ערכי מינימום ומקסימום מוחלטים של הפונקציה הבאה בתחום הגדרתה:
\[ f(x)=\frac{\sqrt{x^2-1}}{x^2} \]
או הסבירו מדוע אינם קיימים.
\end{question}

\begin{hint}
מצאו קודם את תחום ההגדרה, ושימו לב שהפונקציה זוגית ואי-שלילית. בדקו גם את הגבול באינסוף.
\end{hint}

\begin{solution}
\textbf{תחום הגדרה:} $x^2-1\ge0$ ו-$x\neq0$, כלומר $D=(-\infty,-1]\cup[1,\infty)$.
$f$ זוגית, ולכן מספיק לחקור את $[1,\infty)$.

\textbf{מינימום מוחלט:} $f(x)\ge0$ לכל $x\in D$, ו-$f(\pm1)=0$. לכן המינימום המוחלט הוא $\boxed{0}$, המתקבל ב-$x=\pm1$.

\textbf{מקסימום מוחלט:} עבור $x>1$:
\[ f'(x)=\frac{\frac{x}{\sqrt{x^2-1}}\cdot x^2-2x\sqrt{x^2-1}}{x^4}=\frac{x^2-2(x^2-1)}{x^3\sqrt{x^2-1}}=\frac{2-x^2}{x^3\sqrt{x^2-1}}. \]
לכן $f'>0$ ב-$(1,\sqrt2)$ ו-$f'<0$ ב-$(\sqrt2,\infty)$: $f$ עולה ב-$[1,\sqrt2]$ ויורדת ב-$[\sqrt2,\infty)$ (רציפה בקצה $1$). מכאן שלכל $x\ge1$ מתקיים $f(x)\le f(\sqrt2)$, ו-
\[ f(\sqrt2)=\frac{\sqrt{2-1}}{2}=\frac12. \]
(נשים לב גם ש-$\lim_{x\to\infty}f(x)=\lim\frac{\sqrt{1-1/x^2}}{x}=0$, כך שאין "בריחה" באינסוף.) מזוגיות, אותו דבר נכון ב-$(-\infty,-1]$.

\textbf{תשובה:} המקסימום המוחלט הוא $\boxed{\tfrac12}$, המתקבל ב-$x=\pm\sqrt2$; המינימום המוחלט הוא $\boxed{0}$, המתקבל ב-$x=\pm1$.
\end{solution}
